\begin{solapartat}
\punts{0,25} L'angle de reflexió de la llum en una superfície plana és el mateix que l'angle incident.

\punts{0,25} Per tant, el raig que prové dels peus de la Meritxell i es reflecteix al mirall per arribar finalment als seus ulls es reflectirà a la meitat de l'alçada dels ulls de la Meritxell. És a dir
\[ h_\text{reflexió peus} = \frac{h_\text{ulls}}{2} = \frac{1,50}{2} = 0,75\ \mathrm{m}. \]

\punts{0,25} Per la mateixa raó, per veure el punt més alt del seu cap que està 10~cm per sobre dels seus ulls, el raig haurà de reflectir-se a 5~cm per sobre dels seus ulls:
% DUBTE: el numerador de la primera fracció no es llegeix bé a l'original (només s'hi veu "ulls")
\[ h_\text{reflexió cap} = h_\text{ulls} + \frac{h_\text{cap} - h_\text{ulls}}{2} = 1,50 + \frac{0,10}{2} = 1,55\ \mathrm{m} \]
és a dir a 1,55~m del terra.

\punts{0,25} Per tant l'alçada mínima del mirall serà la resta entre aquestes dues alçades:
\[ h_\text{mirall} = h_\text{reflexió cap} - h_\text{reflexió peus} = 1,55 - 0,75 = 0,80\ \mathrm{m}. \]

\punts{0,25} Esquema:
\figura[0.6]{sol_fig1.png}
\end{solapartat}

\begin{solapartat}
\punts{0,25} El raig que determinarà l'alçada màxima d'en Guifré és el que passa pels ulls de la Meritxell i es reflecteix a l'extrem del mirall, a una alçada de 1,55~m.

\punts{0,25} L'angle d'incidència d'aquest feix compleix $\tan\alpha = \dfrac{5\ \mathrm{cm}}{d}$ per tant també l'angle reflectit.

\punts{0,25} Per tant l'alçada d'aquest feix respecte l'extrem del mirall serà $2d\tan\alpha = 10\ \mathrm{cm}$.

\punts{0,25} Per tant l'alçada màxima del Guifré és 10~cm més que l'extrem superior del mirall:
\[ h_\text{màx Guifre} = 1,55\ \mathrm{m} + 10\ \mathrm{cm} = 1,65\ \mathrm{m}. \]

\punts{0,25} Esquema:
\figura[0.65]{sol_fig2.png}
\end{solapartat}
